% !TEX root = ../main.tex

\chapter{相关方法与理论基础}\label{ch:prelim}

本章对信息技术基础理论、系统架构方法、性能优化技术和评估指标体系四个方面的背景知识进行系统梳理。第~\ref{sec:ch2-fundamentals}~介绍信息技术的基本概念和分类，阐述数据处理、系统集成和服务架构的工作原理和适用场景；第~\ref{sec:ch2-optimization}~详细介绍资源调度、负载均衡和弹性伸缩等优化方法，分析不同优化策略的优缺点和适用条件；第~\ref{sec:ch2-performance}~综述性能监控、容错处理和安全管理等技术，并讨论它们在不同场景中的效果和局限性；第~\ref{sec:ch2-evaluation}~介绍系统评估的常用指标和方法，为后续研究提供理论基础和评价标准。

\section{信息技术基础}\label{sec:ch2-fundamentals}

机器学习（Machine Learning, ML）是从数据中自动学习规律和模式的方法，本质上通过``特征提取---模式识别---任务预测/决策''实现从原始数据到任务输出的智能映射。在数据处理场景中，输入特征往往更关注数据的统计特性和结构信息，模型参数则逐渐转向模式识别和决策逻辑，这种从数据到知识的转换过程，使得机器学习成为分类、回归、聚类和生成等任务的核心建模工具。

\subsection{系统架构设计}\label{sec:ch2-system-structure}

从组成上看，典型的信息系统一般由数据层、业务层、服务层和接口层组成。其中，数据层负责存储和管理基础数据；业务层处理核心业务逻辑；服务层提供可复用的功能组件；接口层则对外提供统一的服务接入点。对于最基础的系统架构，其功能映射可表示为：
\begin{equation}
  O = f(I, P)
\end{equation}
其中，\(I\) 表示输入数据，\(P\) 表示系统参数，\(O\) 表示输出结果，是系统的核心要素。若业务为查询型操作，通常返回结构化的数据响应；若业务为事务型操作，则执行相应的业务逻辑并返回处理结果。

在处理复杂业务时，单体架构虽然简单直观，但扩展性和维护性有限。因此，分布式架构如微服务、事件驱动等成为更常见的选择。这些架构通过服务拆分和异步通信来解耦复杂系统。随着业务量的增长，集群部署如负载均衡、故障转移等通过多个节点协同来提升系统的可用性和可靠性。此外，云原生架构通过容器化和DevOps实践，能够实现资源的弹性调度和服务的快速迭代，成为构建现代化系统的主流方法。信息系统的典型架构和各组件的关系如图~\ref{fig:ch2-system-structure}~所示。

\begin{figure}[htbp]
  \centering
  \ThesisFigFillWidth{fig-ch2-system-structure.png}
  \caption{智能系统基础架构与数据流转示意图}
  \label{fig:ch2-system-structure}
\end{figure}

从建设角度看，信息系统的开发主要依赖需求分析、架构设计、技术选型和运维部署四个步骤。给定需求 \(R\) 和约束条件 \(C\)，系统架构 \(S\) 的优化目标可表示为：
\begin{equation}
  \operatorname{minimize} Cost(S) \quad \text{subject to} \quad Performance(S) \geq P_{\min}, \quad Reliability(S) \geq R_{\min}
\end{equation}

其中，\(Cost\) 为建设成本，\(Performance\) 和 \(Reliability\) 为性能和可靠性指标。对于事务处理系统，通常要求高并发和低延迟；对于批处理系统，则更注重吞吐量和数据一致性。由于信息系统在业务支撑中需要平衡功能实现和非功能性需求，同时建设过程高度依赖业务理解和技术选型，因此一旦需求理解存在偏差，或架构设计不当，系统可能在运行过程中出现性能瓶颈或可靠性问题，这正是系统优化需要关注的重要方面。

\section{系统优化方法}\label{sec:ch2-optimization}

模型优化（Model Optimization）旨在通过调整算法和参数来提升模型的性能和效率。按照优化目标可分为准确率优化、效率优化和鲁棒性优化；按照优化方法可分为梯度下降法、进化算法和贝叶斯优化等不同技术。对模型优化的评估通常至少包含三类指标，第一类是模型性能，如准确率、精确率、召回率或F1分数，用于衡量模型的预测能力；第二类是训练效率，如训练时间、收敛速度和资源消耗，用于衡量优化的效率；第三类是泛化能力，如交叉验证性能、在测试集上的表现，用于衡量模型的适用范围。

\subsection{监督学习优化}\label{sec:ch2-supervised}

监督学习是机器学习中最常用的范式之一，其基本目标是学习一个从输入特征到输出标签的映射函数。给定训练数据集 \(D=\{(x_i,y_i)\}\)，监督学习的优化目标是最小化预测输出与真实标签之间的差距。常见的监督学习方法包括线性回归、逻辑回归、支持向量机和神经网络等。

从方法演进来看，最早的监督学习方法主要基于统计学理论，如线性回归和逻辑回归。这些方法虽然简单高效，但在处理复杂非线性问题时表现有限。此后，支持向量机通过核技巧能够处理非线性分类问题，但在大规模数据集上计算成本较高。随着深度学习的发展，神经网络通过多层结构和非线性激活函数，能够自动学习特征表示，在各种任务中取得了优异性能。

在优化策略方面，监督学习主要分为批处理学习和在线学习两种模式。批处理学习一次性使用所有训练数据进行模型训练，适用于数据集规模不大的情况；在线学习则逐个或小批量处理训练数据，能够适应数据流场景，并实时更新模型。近年来，半监督学习和弱监督学习通过利用未标注数据或弱标注信息，进一步提高了学习效率和性能。

\subsection{无监督学习与表示学习}\label{sec:ch2-unsupervised}

无监督学习是从无标注数据中发现隐藏结构和模式的方法，其目标不是预测标签，而是理解数据的内在分布和结构。常见的无监督学习方法包括聚类、降维、密度估计和生成模型等。

聚类算法旨在将相似的数据点分组到一起，常用的方法有K-means、层次聚类和DBSCAN等。这些方法通过定义距离度量或相似性度量，将数据划分到不同的簇中，适用于客户分群、异常检测等场景。

降维技术旨在将高维数据映射到低维空间，同时保留最重要的信息。主成分分析（PCA）通过线性变换寻找方差最大的方向；t-SNE和UMAP等非线性降维方法则保持数据点之间的局部结构，常用于数据可视化和特征提取。

表示学习是从原始数据中学习有意义的表示的方法。自编码器通过编码器-解码器结构学习数据的压缩表示；生成对抗网络（GAN）通过对抗训练学习数据的分布；对比学习则通过对比正负样本学习鲁棒的特征表示。这些方法在无监督预训练中取得了显著成功。

\section{性能提升技术}\label{sec:ch2-performance}

鲁棒性提升（Robustness Enhancement）旨在增强模型对噪声、异常值和分布偏移的抵抗能力。按照技术路线可分为数据层面优化、模型结构优化和训练策略优化等不同方法。对鲁棒性提升的评估通常包含模型在噪声数据、对抗样本和分布外数据上的性能表现。

\subsection{数据层面优化}\label{sec:ch2-data}

数据层面优化主要通过预处理和增强技术来提高数据质量。数据清洗包括处理缺失值、异常值检测和离群点移除等操作，确保训练数据的可靠性和一致性。数据标准化通过均值归一化和方差标准化，使不同特征具有可比性。归一化则将数据缩放到特定范围，如[0,1]或[-1,1]，有助于稳定训练过程。

数据增强通过对现有数据进行各种变换生成新的训练样本，从而增加数据多样性。在图像处理中，常用的增强方法包括旋转、缩放、裁剪、颜色变换和添加噪声等；在文本处理中，可以通过同义词替换、句式变换和回译等技术生成多样化的文本样本。

特征工程是从领域知识出发，设计和选择能有效表示问题的特征。特征选择通过相关性分析、信息增益或L1正则化等方法选择最重要的特征；特征提取通过主成分分析、线性判别分析或流形学习等方法将原始特征映射到新的特征空间；特征构建则通过组合原始特征创建更有表达力的新特征。

\subsection{模型结构优化}\label{sec:ch2-architecture}

模型结构优化通过改变网络设计和添加约束来提高鲁棒性。正则化技术通过在损失函数中添加惩罚项来限制模型复杂度，常见的有L1正则化、L2正则化、Dropout和早停等。这些方法能够有效防止过拟合，提高模型的泛化能力。

集成方法通过组合多个基学习器来提升整体性能。Bagging方法如随机森林通过自助采样训练多个独立模型，然后通过投票或平均得到最终预测；Boosting方法如梯度提升树通过顺序训练多个模型，每个模型专注于纠正前一个模型的错误；Stacking则通过元学习器组合多个基模型的预测结果。

多任务学习通过同时学习多个相关任务来共享知识，提高模型的泛化能力。任务的相似性使得模型能够学习到更通用的特征表示，从而在各个任务上都表现更好。迁移学习则将在源任务上学习到的知识迁移到目标任务，加速目标任务的训练并提高性能。

\subsection{训练策略优化}\label{sec:ch2-training}

训练策略优化通过调整训练过程来提高鲁棒性。优化算法的选择对模型性能有重要影响。随机梯度下降（SGD）虽然简单，但在许多任务中表现出色；带动量的SGD能够加速收敛并跳出局部最优；自适应学习率算法如Adam、RMSprop等能够自动调整学习率，提高训练效率。

对抗训练通过在训练过程中加入对抗样本来提高模型的鲁棒性。对抗样本是对输入进行微小扰动后导致模型错误预测的样本，通过在训练中显式地包含这些样本，模型能够学习到更加鲁棒的特征表示。知识蒸馏则通过训练一个简单的学生模型来模仿复杂教师模型的输出，在保持性能的同时提高模型的鲁棒性。

课程学习通过从简单到复杂的顺序调整训练数据来提高学习效率。初期使用简单或干净的数据，随着训练的进行逐渐引入更复杂或带噪声的数据，使模型逐步适应各种挑战。这种渐进式的训练方式有助于模型学习到更加鲁棒的特征表示。

\section{系统评估方法}\label{sec:ch2-evaluation}

系统评估（System Evaluation）旨在客观地衡量系统在不同场景下的性能表现。按照评估场景可分为功能评估和非功能评估；按照评估方法可分为实验室评估和生产评估。对系统评估需要考虑功能性、可靠性、性能和安全性等多个维度。

\subsection{性能指标}\label{sec:ch2-metrics}

性能指标是量化模型预测能力的数值标准。对于分类任务，准确率是最常用的指标，但可能在不平衡数据上产生误导；精确率和召回率则分别关注预测为正的样本中真正为正的比例，以及所有真正为正的样本中被正确预测的比例；F1分数是精确率和召回率的调和平均，能够平衡两者。对于多分类问题，宏平均和微平均提供了不同的聚合方式。

对于回归任务，常用的指标包括均方误差（MSE）、平均绝对误差（MAE）和R平方等。MSE对大误差惩罚更重，MAE则对所有误差一视同仁；R平方表示模型解释的方差比例，范围在0到1之间，越接近1表示模型拟合越好。对于时序数据，还可以考虑预测值的趋势一致性等指标。

对于排序和推荐任务，常用指标包括准确率@k、召回率@k、平均精度均值（MAP）和归一化折损累计增益（NDCG）等。这些指标考虑了预测结果的排序质量，特别适合信息检索和推荐系统等需要返回有序列表的场景。

\subsection{评估策略}\label{sec:ch2-strategy}

评估策略决定了如何划分和使用数据进行性能测试。简单交叉验证将数据随机分为训练集和测试集；K折交叉验证将数据分成K份，轮流使用其中一份作为测试集，其余作为训练集，取平均性能作为最终结果；分层交叉验证则保证每折中各类样本的比例与原始数据相同。

时间序列评估需要考虑数据的时序特性，通常采用滚动窗口或向前预测的方式。在金融和医疗等高风险领域，还需要进行前瞻性验证，使用未来的数据来评估模型在实际应用中的表现。

鲁棒性评估通过向测试数据添加噪声或扰动来测试模型的稳定性。对抗样本测试评估模型在恶意扰动下的表现；噪声鲁棒性测试评估模型在不同噪声水平下的性能；分布偏移测试则评估模型在不同数据分布上的泛化能力。

\subsection{可解释性评估}\label{sec:ch2-explainability}

可解释性评估衡量模型决策过程的透明度和可理解性。特征重要性分析通过计算特征对预测的贡献来识别关键因素；局部解释如LIME和SHAP能够解释单个预测的依据；全局可视化如注意力热图展示了模型的整体关注模式。

用户评估通过问卷调查等方式评估解释结果的质量和可用性。理解度衡量用户能否理解解释内容；可信度评估用户对解释结果的信任程度；满意度则综合评估用户对整体解释体验的评价。这些指标对于确保解释方法在实际应用中的有效性至关重要。

表~\ref{tab:ch2-supervised-compare}~汇总了各代表性技术方法的特点对比。总体而言，机器学习优化技术的发展主线大致分为三类，第一类，方法从简单的参数优化发展为复杂的结构设计和多任务学习；第二类，技术假设从理想数据条件扩展到实际应用中的噪声和分布偏移场景；第三类，评估指标从单一的准确性发展到综合考虑鲁棒性、效率和可解释性的多维度评价体系。

\begin{table}[htbp]
  \centering
  \footnotesize
  \setlength{\tabcolsep}{4.5pt}%
  \renewcommand{\arraystretch}{1.28}%
  \caption{代表性机器学习方法对比}\label{tab:ch2-supervised-compare}
  \begin{tabularx}{\linewidth}{@{}
    >{\centering\arraybackslash}p{0.168\linewidth}
    >{\centering\arraybackslash}p{0.118\linewidth}
    >{\centering\arraybackslash}p{0.132\linewidth}
    >{\RaggedRight\arraybackslash}X
    >{\RaggedRight\arraybackslash}X
  @{}}
    \toprule\noalign{}
    \textbf{方法} & \textbf{学习类型} & \textbf{适用场景} & \textbf{关键机制} & \textbf{主要特征} \\
    \midrule\noalign{}
    线性回归
      & 回归问题 & 连续数值预测
      & 最小化均方误差，学习线性映射关系
      & 简单高效，可解释性强，但只能处理线性关系 \\
    支持向量机 \cite{svmdoc}
      & 分类问题 & 小到中等规模数据
      & 寻找最优分类超平面，最大化分类间隔
      & 理论基础扎实，通过核技巧处理非线性问题 \\
    决策树
      & 分类/回归 & 特征具有明显层次关系
      & 递归分裂特征空间，构建树形决策结构
      & 直观易理解，无需特征缩放，但容易过拟合 \\
    神经网络
      & 分类/回归/生成等 & 大规模复杂数据
      & 多层非线性映射，自动学习特征表示
      & 表达能力强，适合复杂任务，但需要大量数据和计算资源 \\
    随机森林
      & 分类/回归/特征重要性 & 高维数据
      & 多棵决策树的集成投票
      & 抗过拟合，处理高维数据，提供特征重要性 \\
    梯度提升树
      & 分类/回归 & 结构化数据
      & 顺序训练多个弱学习器，逐步纠正错误
      & 预测精度高，对参数敏感，计算复杂度较高 \\
    自编码器
      & 无监督/表示学习 & 特征提取和降维
      & 通过编码器-解码器结构学习数据的压缩表示
      & 自动学习特征表示，可用于降噪和异常检测 \\
    生成对抗网络
      & 无监督/生成学习 & 数据生成和增强
      & 通过生成器和判别器的对抗训练学习数据分布
      & 能够生成逼真的样本，训练不稳定 \\
    \bottomrule\noalign{}
\end{tabularx}
\end{table}

\section{深度学习基础}\label{sec:ch2-deeplearning}

深度学习（Deep Learning）是机器学习的一个分支，使用多层神经网络来学习数据的层次化表示。从简单的感知机到复杂的变换器模型，深度学习在计算机视觉、自然语言处理和语音识别等领域取得了突破性进展。

\subsection{神经网络基础}\label{sec:ch2-nn}

神经网络是由大量相互连接的处理单元（神经元）组成的计算模型。每个神经元接收输入信号，通过权重进行加权求和，经过激活函数处理后输出。多层神经网络通过堆叠多个隐藏层来学习复杂的特征表示。

前馈神经网络是最基本的网络结构，信号从输入层流向输出层，没有反馈连接。卷积神经网络（CNN）通过卷积层和池化层处理图像数据，利用局部连接和权值共享减少参数数量。循环神经网络（RNN）通过循环连接处理序列数据，具有记忆能力。长短期记忆网络（LSTM）和门控循环单元（GRU）是RNN的改进版本，能够更好地处理长序列。

\subsection{预训练模型}\label{sec:ch2-pretraining}

预训练模型在大规模无标注数据上训练，然后在下游任务上进行微调。这种迁移学习策略能够利用大规模数据学习通用特征，减少对标注数据的依赖。

BERT模型通过掩码语言建模和下一句预测任务预训练，能够理解上下文语义。GPT系列模型使用自回归语言建模预训练，擅长文本生成。ViT将图像分割成固定大小的块，使用Transformer架构处理图像数据，在视觉任务中表现出色。

预训练模型的微调策略包括适配头添加、参数解冻和提示学习等。适配头在预训练模型上添加任务特定的输出层；参数解冻允许微调整个模型的参数；提示学习则通过输入模板引导模型生成特定格式的输出。

\subsection{注意力机制}\label{sec:ch2-attention}

注意力机制模拟人类选择性关注的能力，使模型能够聚焦于输入中的重要部分。自注意力计算序列中所有元素之间的依赖关系，通过查询（Query）、键（Key）和值（Value）的交互得到注意力权重。

多头注意力将注意力分成多个"头"，每个头关注不同的表示子空间。交叉注意力用于处理不同模态或不同序列之间的交互。层次化注意力在多个粒度上应用注意力，从局部到全局逐步聚合信息。

注意力机制的可视化能够解释模型的决策过程，通过热图显示模型关注哪些输入部分。这种透明性对于建立用户信任和调试模型具有重要意义。